\documentclass[11pt,a4paper]{article}
\usepackage[T1]{fontenc}
\usepackage[utf8]{inputenc}
\usepackage{amsmath,amssymb,xcolor}
\usepackage[margin=25mm]{geometry}
\usepackage[colorlinks=true,urlcolor=blue]{hyperref}
\definecolor{nexorablue}{HTML}{2459B5}
\setlength{\parindent}{0pt}
\setlength{\parskip}{7pt}
\setlength{\emergencystretch}{2em}
\renewcommand{\familydefault}{\sfdefault}
\begin{document}

{\small\bfseries\color{nexorablue} NEXORA CPGE}\par
{\small MP / WKB-SPE-SI-02 / Énoncé}\par
{\LARGE\bfseries Réponse impulsionnelle et réponse indicielle}\par
\vspace{0.25em}\hrule\vspace{0.6em}

Pour un système de fonction de transfert $H(p)$ et initialement au repos, exprimer la sortie dans le domaine de Laplace lorsque l’entrée est une impulsion de Dirac, puis lorsqu’elle est un échelon unitaire. Quelle relation temporelle relie les deux réponses causales ?

\bigskip
{\footnotesize Adaptation française par Nexora CPGE. Source : Contributeurs de Wikibooks,
\href{https://en.wikibooks.org/wiki/Control_Systems/Transfer_Functions\#Using_the_Transfer_Function}{\textit{Control Systems}}, Transfer Functions, exemples Impulse Response et Step Response.
\href{https://creativecommons.org/licenses/by-sa/4.0/}{CC BY-SA 4.0}.}
\end{document}
